\documentclass[aspectratio=169,10pt]{beamer}

% ═════════════════════════════════════════════════════════════════════════════
%  TEMPLATE LIBRARY DEMO
%  Showcases all 4 themes × 8 layouts × 3 component categories
% ═════════════════════════════════════════════════════════════════════════════

\usepackage[academic]{template-lib/template-lib}
\uselayout{text}
\uselayout{1img}
\uselayout{2img}
\uselayout{eq}
\uselayout{twocol}
\uselayout{table}
\uselayout{3img}
\uselayout{imgtop}

\TLsetfoot{Template Library Demo · 2026}

\title{Template Library}
\subtitle{Three-Tier Beamer System}
\author{Demo Document}
\institute{Theme + Layout + Component}
\date{2026年5月23日}

\begin{document}

% ── Title Slide ─────────────────────────────────────────────────────────────
\TLtitlecenter
  {Template Library}
  {Theme · Layout · Component}
  {Demo Document}
  {2026年5月23日}

% ── S1: Overview ────────────────────────────────────────────────────────────
\TLsection{System Overview}

\begin{frame}{Three-Tier Architecture}
  \begin{TLtext}
    \textbf{Tier 1 — Themes:} Color palettes + typography\\
    \quad academic · teal · dark · navygold
    
    \vspace{0.8em}
    \textbf{Tier 2 — Layouts:} Page structure templates\\
    \quad text · 1img · 2img · 3img · eq · table · imgtop · twocol
    
    \vspace{0.8em}
    \textbf{Tier 3 — Components:} Reusable content blocks\\
    \quad title · block · fig
  \end{TLtext}
  
  \TLtakeaway{%
    Mix and match any theme with any layout. Components work everywhere.%
  }
\end{frame}

% ── S2: Theme Showcase ──────────────────────────────────────────────────────
\TLsection{Theme Tier}

\begin{frame}{Theme: Academic (Navy + Brick)}
  \begin{TLtext}
    \begin{itemize}
      \item Primary: \TLterm{Deep Navy} \textcolor{TLprimary}{\rule{1cm}{0.5em}}
      \item Accent: \TLterm{Brick Red} \textcolor{TLaccent}{\rule{1cm}{0.5em}}
      \item Best for: General academic presentations
    \end{itemize}
  \end{TLtext}
  
  \TLinfoblock{Info Block Example}{%
    This is how an information block looks in the academic theme.
    It uses the primary color for the border and tint for background.%
  }
  
  \vspace{0.3em}
  \TLalertblock[Alert]{%
    This is an alert block using the accent color.%
  }
\end{frame}

\begin{frame}{Theme Comparison}
  \TLtwocol{
    \textbf{Academic}\\
    \textcolor{TLprimary}{\rule{1cm}{0.5em}} Navy\\
    \textcolor{TLaccent}{\rule{1cm}{0.5em}} Brick
    
    \vspace{0.5em}
    \textbf{Teal}\\
    \textcolor{TLprimary}{\rule{1cm}{0.5em}} Teal\\
    \textcolor{TLaccent}{\rule{1cm}{0.5em}} Amber
  }{
    \textbf{Dark}\\
    \textcolor{TLprimary}{\rule{1cm}{0.5em}} Soft Blue\\
    \textcolor{TLaccent}{\rule{1cm}{0.5em}} Gold
    
    \vspace{0.5em}
    \textbf{NavyGold}\\
    \textcolor{TLprimary}{\rule{1cm}{0.5em}} Navy\\
    \textcolor{TLaccent}{\rule{1cm}{0.5em}} Gold
  }
\end{frame}

% ── S3: Layout Showcase ─────────────────────────────────────────────────────
\TLsection{Layout Tier}

\begin{frame}{Layout: Text Only}
  \begin{TLtext}
    \textbf{Pure text flow with takeaway support}
    
    \begin{enumerate}
      \item First point with \TLterm{highlighted term}
      \item Second point with \TLyes\ positive marker
      \item Third point with \TLno\ negative marker
    \end{enumerate}
  \end{TLtext}
  
  \TLtakeaway{%
    The \texttt{text} layout is the simplest — just content flowing naturally.%
  }
\end{frame}

\begin{frame}{Layout: Single Image}
  \TLoneimgleft[0.50]{sparse-view-ct-1.jpg}{
    \textbf{CT Imaging Process}
    \begin{itemize}
      \item X-ray projection
      \item Sinogram recording
      \item FBP reconstruction
    \end{itemize}
    
    \vspace{0.5em}
    \TLresultblock{Result}{PSNR +2.3 dB improvement}
  }
\end{frame}

\begin{frame}{Layout: Two Images}
  \TLtwoimg[0.45]
    {ours_recon_18v.png}{18-view (severe)}
    {ours_recon_72v.png}{72-view (mild)}
  
  \vspace{0.5em}
  \begin{minipage}{\textwidth}
    \centering
    \TLtakeaway{%
      View count directly affects reconstruction quality.%
    }
  \end{minipage}
\end{frame}

\begin{frame}{Layout: Equation Focus}
  \TLeqsingle{
    \mathcal{A}f(\theta,s)=\int_{-\infty}^{+\infty}\!
    f(s\cos\theta - t\sin\theta,\;
      s\sin\theta + t\cos\theta)\,dt
  }{%
    The Radon transform maps an image to its sinogram representation.
    Each point $(\theta,s)$ is a line integral along the corresponding ray.%
  }
\end{frame}

\begin{frame}{Layout: Two-Column Comparison}
  \TLtwocol{
    \textbf{Forward Problem}
    \begin{itemize}
      \item Given: Image $f(x,y)$
      \item Compute: Projections
      \item Operator: $\mathcal{A}$
    \end{itemize}
  }{
    \textbf{Inverse Problem}
    \begin{itemize}
      \item Given: Sinogram
      \item Compute: Image
      \item Operator: $\mathcal{A}^\dagger$
    \end{itemize}
  }
  
  \vspace{0.5em}
  \TLtwocoldiv{
    \centering
    \textbf{Analytical}\\
    FBP, Filtered Backprojection
  }{
    \centering
    \textbf{Iterative}\\
    SART, SIRT, CG
  }
\end{frame}

% ── S4: Component Showcase ──────────────────────────────────────────────────
\TLsection{Component Tier}

\begin{frame}{Block Components}
  \TLinfoblock{Information Block}{%
    Use for definitions, explanations, or background info.%
  }
  
  \vspace{0.3em}
  \TLalertblock[Important Note]{%
    Use for key insights, warnings, or critical takeaways.%
  }
  
  \vspace{0.3em}
  \TLresultblock{Experimental Result}{%
    Use for positive results, improvements, or achievements.%
  }
  
  \vspace{0.3em}
  \TLwarnblock{Limitation}{%
    Use for known issues, constraints, or future work.%
  }
\end{frame}

\begin{frame}{Figure Components}
  \TLtwocol{
    \TLautoimg[max height=0.35\textheight]{sparse-view-ct-1.jpg}
    
    \vspace{0.3em}
    \centering
    \footnotesize Auto-scaled image
  }{
    \TLimgcap[max height=0.35\textheight]{sparse-view-ct-2.jpg}{%
      Image with automatic caption%
    }
  }
  
  \vspace{0.5em}
  \TLtakeaway{%
    \texttt{TLautoimg} fits to column width.
    \texttt{TLimgcap} adds caption below.%
  }
\end{frame}

\begin{frame}{Data Table}
  \TLtablefull{
    \begin{tabular}{lccc}
      \toprule
      \TLthead Method & 18-view & 36-view & 72-view \\
      \midrule
      FBP & 24.5 & 28.3 & 32.1 \\
      \TLthl U-Net & 31.2 & 34.5 & 36.8 \\
      Ours & \textbf{35.8} & \textbf{38.2} & \textbf{39.5} \\
      \bottomrule
    \end{tabular}
  }{PSNR (dB) comparison on AAPM-LDCT dataset}
\end{frame}

% ── S5: Combined Example ────────────────────────────────────────────────────
\TLsection{Real-World Example}

\begin{frame}{CvG-Diff: Complete Slide}
  \TLoneimgright[0.42]{
    \textbf{Cross-view Generalized Diffusion}
    \begin{itemize}
      \item Single model for all view counts
      \item Cold Diffusion framework
      \item EPCT training strategy
    \end{itemize}
    
    \vspace{0.3em}
    \TLresultblock{Key Result}{%
      18-view: \TLdeltaup{4.18 dB} vs baseline%
    }
  }{paper_fancy_title.jpg}
\end{frame}

\begin{frame}{Method Comparison}
  \TLprocon{
    \begin{itemize}
      \item[\TLyes] High quality
      \item[\TLyes] Cross-view generalization
      \item[\TLyes] Low NFE (50 steps)
    \end{itemize}
  }{
    \begin{itemize}
      \item[\TLno] Training time
      \item[\TLno] Memory intensive
      \item[\TLno] Hyperparameter tuning
    \end{itemize}
  }
\end{frame}

% ── End ─────────────────────────────────────────────────────────────────────
\TLsection{Thank You}

\begin{frame}[plain]
  \vfill
  \centering
  {\Huge\bfseries\color{TLprimary} Thank You}
  
  \vspace{1em}
  {\large Questions?}
  
  \vspace{2em}
  {\small\color{TLinkSoft}
    Template Library — Three-Tier Beamer System\\
    \texttt{github.com/yourname/template-lib}
  }
  \vfill
\end{frame}

\end{document}
